\subsection*{Problema 2 (b)}
\emph{Solución.}
Al añadir la columna \(n\):
\begin{itemize}
\item Si no hay torre en la columna \(n\), corresponde a \(T(n-1,k-1)\) formas.
\item Si hay torre en la columna \(n\), hay \(k\) filas disponibles y \(T(n-1,k)\) formas previas, luego \(k\,T(n-1,k)\).
\end{itemize}
Así,
\[
T(n,k) = T(n-1,k-1) + k\,T(n-1,k),
\]
con \(T(0,0)=1\) y \(T(n,k)=0\) si \(k>n\) o \(k<0\). Es la recurrencia de los números de Stirling:
\[
T(n,k)=\left\{n\atop k\right\}.
\]
